\documentclass[12pt,a4paper]{article}
\usepackage[UTF8]{ctex}
\usepackage{amsmath,amssymb,amsthm}
\usepackage{geometry}

\geometry{margin=2.5cm}

\title{SE(3)群上的对数映射：$[X_e] = \log(X^{-1} X_d)$的理解}
\author{第11章补充说明}
\date{}

\begin{document}

\maketitle

\section{问题}

\textbf{如何理解$[X_e] = \log(X^{-1} X_d)$中的$\log$？}

这个公式出现在任务空间控制中，用于计算SE(3)群上的误差。

\section{SE(3)群回顾}

\subsection{定义}

SE(3)是特殊欧氏群，元素是4×4齐次变换矩阵：
\begin{equation}
X = \begin{bmatrix}
R & p \\
0 & 1
\end{bmatrix} \in SE(3)
\end{equation}
其中：
\begin{itemize}
\item $R \in SO(3)$是3×3旋转矩阵
\item $p \in \mathbb{R}^3$是平移向量
\end{itemize}

\subsection{SE(3)的性质}

\begin{itemize}
\item 群乘法：$X_1 X_2 = \begin{bmatrix} R_1 R_2 & R_1 p_2 + p_1 \\ 0 & 1 \end{bmatrix}$
\item 逆：$X^{-1} = \begin{bmatrix} R^T & -R^T p \\ 0 & 1 \end{bmatrix}$
\item 单位元：$I_4 = \begin{bmatrix} I_3 & 0 \\ 0 & 1 \end{bmatrix}$
\end{itemize}

\section{李代数se(3)}

\subsection{定义}

se(3)是SE(3)的李代数，元素是4×4反对称矩阵（twist）：
\begin{equation}
[\xi] = \begin{bmatrix}
[\omega] & v \\
0 & 0
\end{bmatrix} \in se(3)
\end{equation}
其中：
\begin{itemize}
\item $[\omega] \in so(3)$是3×3反对称矩阵，对应角速度向量$\omega \in \mathbb{R}^3$
\item $v \in \mathbb{R}^3$是线速度向量
\end{itemize}

\subsection{反对称矩阵}

对于向量$\omega = [\omega_1, \omega_2, \omega_3]^T$，反对称矩阵为：
\begin{equation}
[\omega] = \begin{bmatrix}
0 & -\omega_3 & \omega_2 \\
\omega_3 & 0 & -\omega_1 \\
-\omega_2 & \omega_1 & 0
\end{bmatrix}
\end{equation}

\section{对数映射$\log: SE(3) \to se(3)$}

\subsection{定义}

对于$X = \begin{bmatrix} R & p \\ 0 & 1 \end{bmatrix} \in SE(3)$，对数映射$\log(X)$计算对应的twist $[\xi]$。

\subsection{计算步骤}

\textbf{步骤1：提取旋转和平移}
\begin{equation}
R = X[1:3, 1:3], \quad p = X[1:3, 4]
\end{equation}

\textbf{步骤2：计算旋转的对数映射}

对于旋转矩阵$R$，使用SO(3)的对数映射：
\begin{equation}
[\omega] = \log(R) \in so(3)
\end{equation}

SO(3)的对数映射公式：
\begin{align}
\theta &= \arccos\left(\frac{\text{tr}(R) - 1}{2}\right) \\
[\omega] &= \begin{cases}
\frac{\theta}{2\sin\theta}(R - R^T) & \text{如果} \theta \neq 0 \\
0 & \text{如果} \theta = 0
\end{cases}
\end{align}

从$[\omega]$提取角速度向量$\omega$（反对称矩阵的"vee"操作）。

\textbf{步骤3：计算平移部分}

\begin{equation}
v = \left(\frac{1}{\theta^2}\left(I - \frac{[\omega]}{2}\right) - \frac{[\omega]^2}{2\theta\sin\theta}\right) p
\end{equation}

当$\theta \to 0$时（小角度近似）：
\begin{equation}
v \approx p
\end{equation}

\textbf{步骤4：组合成twist}

\begin{equation}
[\xi] = \begin{bmatrix}
[\omega] & v \\
0 & 0
\end{bmatrix}
\end{equation}

或者表示为6维向量：
\begin{equation}
\xi = \begin{bmatrix} \omega \\ v \end{bmatrix} \in \mathbb{R}^6
\end{equation}

\section{误差表示$[X_e] = \log(X^{-1} X_d)$}

\subsection{物理意义}

\begin{itemize}
\item $X$：实际末端执行器位姿
\item $X_d$：期望末端执行器位姿
\item $X^{-1} X_d$：从实际位姿到期望位姿的相对变换
\item $\log(X^{-1} X_d)$：将这个相对变换映射到李代数se(3)
\item $[X_e]$：误差的twist表示
\end{itemize}

\subsection{为什么使用$X^{-1} X_d$而不是$X_d - X$？}

\begin{enumerate}
\item \textbf{SE(3)不是向量空间：}不能直接做减法
\item \textbf{群结构：}使用群乘法$X^{-1} X_d$表示相对变换
\item \textbf{对数映射：}将群元素映射到向量空间（李代数），可以线性化
\end{enumerate}

\subsection{误差的6维向量表示}

通常将twist表示为6维向量：
\begin{equation}
X_e = \begin{bmatrix} \omega_e \\ v_e \end{bmatrix} \in \mathbb{R}^6
\end{equation}
其中：
\begin{itemize}
\item $\omega_e \in \mathbb{R}^3$：旋转误差（角速度）
\item $v_e \in \mathbb{R}^3$：平移误差（线速度）
\end{itemize}

\section{小角度近似}

当误差很小时，可以简化计算：

\textbf{小角度近似：}
\begin{equation}
X^{-1} X_d \approx I + [X_e]
\end{equation}

其中$[X_e]$是小的twist。

在这种情况下：
\begin{equation}
X_e \approx \begin{bmatrix}
\text{vee}(\log(R^T R_d)) \\
R^T(p_d - p)
\end{bmatrix}
\end{equation}

\section{在控制中的应用}

在任务空间控制中，误差$X_e$用于计算控制律：

\begin{equation}
\dot{\theta} = J^{-1}(\theta)\left(\dot{x}_d + K_p X_e\right)
\end{equation}

其中$X_e$是6维误差向量。

\section{总结}

\textbf{关键点：}

\begin{enumerate}
\item $\log$是李群SE(3)到李代数se(3)的映射
\item $X^{-1} X_d$表示从实际位姿到期望位姿的相对变换
\item $[X_e] = \log(X^{-1} X_d)$将相对变换映射为twist（误差表示）
\item 这个误差可以表示为6维向量$X_e = [\omega_e^T, v_e^T]^T$
\item 在控制中，这个误差用于反馈控制律
\end{enumerate}

\textbf{与标量情况的类比：}

在标量情况下，误差是$e = x_d - x$（直接减法）。

在SE(3)群中，不能直接减法，所以使用：
\begin{equation}
X_e = \log(X^{-1} X_d)
\end{equation}

这相当于"群上的减法"。

\end{document}

